%% Milestone 2 revised proposal, group 38. Build: latexmk -pdf proposal.tex
%% Every number comes from numbers.tex, which ms2.ipynb writes.
\documentclass[11pt]{article}
\usepackage[margin=0.7in]{geometry}
\usepackage[T1]{fontenc}
\usepackage{lmodern}
\usepackage{microtype}
\usepackage[numbers,sort&compress]{natbib}
\usepackage[hidelinks]{hyperref}
\usepackage{enumitem}
\usepackage{titlesec}
\titlespacing*{\section}{0pt}{0.6em}{0.2em}
\titleformat*{\section}{\normalsize\bfseries}
\setlist{nosep, leftmargin=1.3em}
\setlength{\parskip}{0.3em}
\setlength{\parindent}{0pt}
\setlength{\bibsep}{0pt}
% Written by ms2.ipynb; do not edit.
\newcommand{\nSnapshots}{2,720}
\newcommand{\nUniverse}{950}
\newcommand{\nOldTickers}{977}
\newcommand{\nYahoo}{677}
\newcommand{\nAlpaca}{165}
\newcommand{\nNoData}{108}
\newcommand{\nRenames}{30}
\newcommand{\nWrong}{9}
\newcommand{\nBreaks}{4}
\newcommand{\nRawRows}{3,423,011}
\newcommand{\nCleanRows}{2,315,615}
\newcommand{\dataMB}{442}
\newcommand{\nTableRows}{473,684}
\newcommand{\nTableCols}{183}
\newcommand{\nCompanies}{725}
\newcommand{\nWeeks}{1,139}
\newcommand{\covEarly}{55\%}
\newcommand{\covEarlyRaw}{53\%}
\newcommand{\covTest}{99.4\%}
\newcommand{\covEnd}{99.8\%}
\newcommand{\stressShare}{7.3\%}
\newcommand{\nStressWeeks}{89}
\newcommand{\nZeroRange}{80}
\newcommand{\rvCorr}{0.95}
\newcommand{\rvN}{12,515}
\newcommand{\rvGap}{-0.06}


\begin{document}

{\large\textbf{Forecasting Next-Week Volatility of S\&P 500 Stocks: Revised Proposal}}\\[0.1em]
Group 38, project 12: Ashley Lee, Juntang Wang, Shuo Yang, Simon Manivong, Terry Su. TF: Daniel More Torres.
Code and notebook: \url{https://github.com/qqgjyx/AC-209A_VolForecast}

%% =================================================================
%% Question
%% =================================================================
\section*{1. Project question}

For each S\&P 500 stock and week, we forecast next week's log realized variance, measured by the
Garman--Klass estimator from daily prices \citep{garman1980gk}. Models train on 2005--2015, tune on
2016--2018 and are tested on 2019--2026 with yearly expanding refits, a one-week gap at each
boundary and no tuning on test data.
\begin{enumerate}
  \item \textbf{Accuracy.} Do pooled models fit on all stocks (ridge and lasso, then gradient-boosted
        trees) beat a HAR regression fit per stock \citep{corsi2009har} on mean QLIKE
        \citep{patton2011qlike} and MSE? The Diebold--Mariano test \citep{diebold1995dm} runs at 5\% on the
        weekly cross-sectional mean loss difference with Newey--West errors. Per-stock HAR needs 52 weeks of
        history (younger stocks use a pooled HAR), and every model is scored on the same rows.
  \item \textbf{Attribution and regimes.} Which feature blocks carry any gain (refitting without each
        block), and does the gain keep its sign in calm weeks and in stress weeks (VIX $\ge 30$ on the Friday)?
\end{enumerate}
It is answered when, for each model, the test-period table of losses and Diebold--Mariano statistics
against HAR, and the block-ablation table split by regime, exist.

\textbf{How and why it changed.} The economic-value question (volatility-targeted portfolio, Sharpe
ratio, drawdown, Fama--French regression) is dropped, as the staff notes require. Earnings timing comes
from the filing times of SEC 8-K Item 2.02 releases, a proxy for the announcement, not quarterly
filings that arrive weeks after it. Fundamentals from XBRL are deferred: they start in mid-2009 and the question does not
need them. Stress is defined by VIX at prediction time, not by calendar years, and qlib's Alpha158
price-volume factors \citep{yang2020qlib} join the features.

%% =================================================================
%% Data
%% =================================================================
\section*{2. Data}

Four sources from the original proposal: S\&P 500 membership by date \citep{aultman2026sp500}
(\nSnapshots{} snapshots, \nOldTickers{} tickers since 2005, \nUniverse{} companies once renames are
merged); market state from FRED \citep{fred2026} (VIX, 3-month VIX, credit spread, curve slope,
T-bill); daily prices from Yahoo Finance \citep{aroussi2026yfinance} (\nYahoo{} tickers, 2005--2026);
earnings-release filing times and industry codes from SEC EDGAR \citep{sec2026edgar}. Two added since:
Alpaca \citep{alpaca2026data}, whose daily bars from 2016 fill \nAlpaca{} members Yahoo lacks and whose
5-minute bars check the target, and CRSP \citep{crsp2026}, every member's daily prices to 2025 through
WRDS, for the survivorship check. The raw table has one row per stock and trading day, \nRawRows{} rows
(\dataMB{} MB), \nCleanRows{} after cleaning. The modeling table has \textbf{\nTableRows{} rows and
\nTableCols{} columns}: one row per stock and week (\nCompanies{} companies, \nWeeks{} weeks), holding
the target and every feature known at that Friday's close.

%% =================================================================
%% Findings
%% =================================================================
\section*{3. What we learned from the data}

\begin{itemize}
  \item \textbf{Survivorship.} Yahoo drops delisted companies. We recovered \nRenames{} renamed tickers
        (BK became BNY) and filled members Yahoo lacks from Alpaca from 2016; \nNoData{} tickers, all gone
        before 2016, have no free data. Coverage of index members is \covEarly{} in 2005, \covTest{} in
        2019 and \covEnd{} in 2026, so pre-2016 training data leans toward survivors. A check against
        CRSP (WRDS, every member) measures the cost: from 2005 to 2015 the free table has
        \crspWeeksFree{} stock-weeks against \crspWeeksCrsp{}, and its mean log variance is \crspGap{}
        below CRSP's in the same year, on average (about 4\% less variance); from 2016 the gap is
        under 0.01, and from 2019 under 1\% of stock-weeks are missing. We propose to train on CRSP
        to 2025.
  \item \textbf{Ticker reuse and wrong series.} Yahoo reuses symbols (today's \texttt{FB} is an ETF), so
        prices are kept only near each ticker's index spells; \nWrong{} series inside a spell belong to
        other securities (Yahoo's \texttt{PARA} trades near \$100{,}000), and \nBreaks{} change company on a date
        (Yahoo's \texttt{JCI} before 2016 is Tyco); hand-checked tables remove both.
  \item \textbf{The target is a sound proxy.} Against 5-minute realized variance for \rvCompanies{} companies
        since 2016, delisted ones included (\rvN{} stock-weeks), log weekly Garman--Klass variance has
        correlation \rvCorr{} with it, and \rvCorrWithin{} within company, easing the proposal's worry
        about noisy daily targets; weeks before 2016 are not covered by this check.
  \item \textbf{Stress is rare.} \stressShare{} of stock-weeks are stress weeks, so results must be
        reported per regime, not as one pooled loss.
  \item \textbf{Leakage.} Feature windows end at the Friday close (unit tests truncate the data and
        recompute), an 8-K filed after 16:00 counts for the next week, membership is point-in-time (the last
        snapshot, 2026-08-18, is carried forward), and
        every price feature is a ratio, so back-adjusted prices cannot leak.
\end{itemize}

%% =================================================================
%% Plan
%% =================================================================
\section*{4. Plan for the rest of the project}

\textbf{Features}: HAR lags (day, week, four weeks); Parkinson, Rogers--Satchell, overnight and downside
variance; returns; dollar volume and Amihud illiquidity; SPY variance; VIX, 3-month VIX, credit spread,
curve slope, T-bill; weeks since earnings, an expected-earnings flag and industry; Alpha158.
\textbf{Models}: per-stock HAR and GARCH(1,1) baselines; pooled ridge and lasso on standardized
features with training-window imputation; gradient-boosted trees (LightGBM). Clustering and neural
networks only if time remains \citep{christensen2023ml,gu2020ml}. \textbf{Evaluation}: QLIKE and MSE on
2019--2026, Diebold--Mariano tests against HAR, block ablation, every result split into calm and
stress weeks.

%% =================================================================
%% Staff notes
%% =================================================================
\section*{5. Staff notes and the AC209A requirement}

Economic-value backtest: dropped. Regularized regression and boosted trees: the two pooled models of
Q1; clustering and neural networks wait. yfinance failures and missing delisted stocks: fetches are
cached (the notebook runs with no keys from a team copy), renames, wrong series and the Alpaca fill are documented, and coverage is
reported. XBRL only from mid-2009: fundamentals deferred. AC209A outside method: HAR, GARCH(1,1), the QLIKE
loss and the Diebold--Mariano test, none covered in class.

\textbf{Use of generative AI.} Claude Code (Anthropic) helped write the data pipeline, tests and this
text; the group reviewed every step and verified the numbers, which the notebook computes and exports.

{\footnotesize
\bibliographystyle{plainnat}
\bibliography{ref}
}

\end{document}
